% ===================================================================
%  図表第 13 号 — ホテル。— 食堂。— カフェ。
%
%  Ceci n'est PAS une transcription : c'est une traduction, faite en
%  2026, en japonais d'aujourd'hui. Voir texto/ja/10-zuhyo-01.tex
%  pour la consigne, qui vaut pour les seize fichiers.
%
%  LES JEUX DE LA CAFETERIA NE SE TRADUISENT PAS PAR LEURS COUSINS
%  JAPONAIS. « shako » est l'echecs d'Europe, non le 将棋 : les deux
%  jeux n'ont ni le meme damier ni les memes pieces, et la gravure
%  montre l'echiquier europeen. On ecrit donc チェス, チェッカー,
%  バックギャモン — en katakana, comme le japonais nomme ces jeux-la.
% ===================================================================

%%K t13-noto-1 noto
\VUnotes{143.2mm}{%
\textit{\VUgras{(*) 部屋の細目については図表第 6 号を見よ。}}%
}

%%K t13-apar-1 apar
\VUsaut{3.0mm}
\VUtitre{15.0pt}{60}{図表第 13 号}
\VUsaut{1.6mm}
\VUpk{10.2pt}{40}{ホテル。--- 食堂。}
\VUsaut{2.0mm}
\VUpk{10.2pt}{40}{カフェ。}
\VUsaut{3.4mm}

%%K t13-01-1 p
1. --- 昨日の夕方ロワイヤンに着いて、私は友人にすすめられた\textit{オセアン・ホテル}
に泊まった。

%%K t13-01-2 p
よい\VUgras{部屋} \textsuperscript{(2)} をあてがわれた。三\VUgras{階}
\textsuperscript{(1)} で、そこで私はぐっすり眠った(*)。

%%K t13-02-1 p
2. --- 今朝、部屋を出た。\VUgras{女中} \textsuperscript{(3)} がそこへ朝食を運んで
くれていたのである。私は\VUgras{喫煙室} \textsuperscript{(5)} に入った。そこは
\VUgras{読書室} \textsuperscript{(4)} も兼ねていて、\VUgras{新聞} \textsuperscript{(6)}
を読みながら\VUgras{巻煙草} \textsuperscript{(8)} を吸った。読み終えると
\VUgras{エレベーター} \textsuperscript{(9)} で下におり、町を散歩しに出かけた。

%%K t13-03-1 p
3. --- ホテルの\VUgras{会計係} \textsuperscript{(12)} に\VUgras{共同食卓}
\textsuperscript{(11)} の食事は何時からかと尋ねるのを忘れていたので、私は早めに戻り、
ついでにホテルの\VUgras{浴室} \textsuperscript{(10)} で湯を浴びた。それから
\VUgras{帳場} \textsuperscript{(15)} へ行って友人の一人に電話をかけようとした。ところが
\VUgras{電話} \textsuperscript{(13)} はふさがっていた。困っている私を見て、
\VUgras{女会計} \textsuperscript{(14)} は——\VUgras{勘定書} \textsuperscript{(17)} を
\VUgras{宿泊者名簿} \textsuperscript{(16)} に書き入れているところだった——
\VUgras{門番} \textsuperscript{(18)} に頼んで\VUgras{ボーイ} \textsuperscript{(19)} を
やり、\VUgras{自転車} \textsuperscript{(20)} で私の用を足させてくれた。

%%K t13-04-1 p
4. --- 私は礼を言って外へ出、\VUgras{赤帽} \textsuperscript{(24)}（雑用係である）
に心づけをやった。彼は\VUgras{手押し車} \textsuperscript{(25)} で停車場から私の
\VUgras{帽子箱} \textsuperscript{(26)}、\VUgras{旅行鞄} \textsuperscript{(27)}、
\VUgras{手提袋} \textsuperscript{(28)} を運んできてくれたのだ。ちょうどそこへ
\VUgras{郵便配達} \textsuperscript{(21)} が来て、一通の\VUgras{手紙}
\textsuperscript{(22)} を渡してくれた。

%%K t13-05-1 p
5. --- 私はなお少しのあいだ戸口に、\VUgras{郵便受け} \textsuperscript{(23)} の
そばに立って、往来のさまを眺めていた。\VUgras{車掌} \textsuperscript{(30)} は——
\VUgras{乗合馬車} \textsuperscript{(29)} の車掌である——\VUgras{大型旅行箱}
\textsuperscript{(31)} を自分の車に積みこんでいた。それは\VUgras{旅客}
\textsuperscript{(32)} のもので、\VUgras{ホテルの主人} \textsuperscript{(33)} がその客に
別れを告げていた。若い\VUgras{電報配達} \textsuperscript{(35)} は少しも急がず、ホテルの
客あての\VUgras{電報} \textsuperscript{(34)} を届けにきた。\VUgras{観光客}
\textsuperscript{(36)} が一人、\VUgras{写真機} \textsuperscript{(38)} を肩から斜めに
かけて、\VUgras{案内書} \textsuperscript{(37)} を調べていた。

%%K t13-tit-1 sub
\VUsaut{4.0mm}
\VUpk{10.2pt}{40}{昼食どき。}
\VUsaut{3.0mm}

%%K t13-06-1 p
6. --- 柵の前では、のんきな\VUgras{使い走り} \textsuperscript{(39)} が
\VUgras{荷担ぎ鉤} \textsuperscript{(40)} と\VUgras{靴磨き箱} \textsuperscript{(41)}
のあいだで居眠りをしていた。若い\VUgras{洗濯女} \textsuperscript{(42)} は洗濯物の
\VUgras{籠} \textsuperscript{(43)} をさげて、戸口に立つ\VUgras{給仕頭}
\textsuperscript{(44)} のいかめしい顔つきを笑っていた。

%%K t13-07-1 p
7. --- \VUgras{料理人} \textsuperscript{(45)} が\VUgras{台所}
\textsuperscript{(47)} ——\VUgras{地下} \textsuperscript{(48)} にある——から出てきて
\VUgras{下働き} \textsuperscript{(46)} を使いにやるのを見て、私は昼食の時刻が近いことを
思い出し、中庭を通って食堂へ入った。中庭では\VUgras{自動車の運転手}
\textsuperscript{(50)} が\VUgras{自動車} \textsuperscript{(49)} を停めており、
\VUgras{機械工} \textsuperscript{(52)} が\VUgras{自転車乗り} \textsuperscript{(53)} の
指図どおりに、事故を起こしたばかりの\VUgras{石油三輪車} \textsuperscript{(51)} を
直していた。

%%K t13-08-1 p
8. --- 若い\VUgras{ボーイ} \textsuperscript{(54)} が\VUgras{ビールのコップ}
\textsuperscript{(55)} を運んでいたので、共同食卓はベランダの下かと尋ねた。そうだと言う
ので、私は一つの卓に席を取り、たいへんよい食事を出させた。

%%K t13-noto-2 noto
\VUnotes{143.2mm}{%
\textit{\VUgras{(*) 語彙に載せた料理の一覧を使えば、教師は次の献立を容易にいろいろ
に変えることができる。}}%
}

%%K t13-tit-2 sub
\VUsaut{4.0mm}
\VUpk{10.2pt}{40}{すばらしい昼食。}
\VUsaut{3.0mm}

%%K t13-09-1 p
9. --- 給仕が昼食の献立(*)と酒の一覧を持ってきた。私は選んで注文した。
\VUgras{前菜には小海老}を\VUgras{バター}添えで、一の皿には\VUgras{カーン風もつ煮}を。
\VUgras{魚}は取らず、羊の\VUgras{もも肉}を一人前、\VUgras{野菜の皿}にはクリームを
添えた\VUgras{ほうれん草}を頼んだ。そのあと\VUgras{甘味}はどれも気に入らなかったので、
いろいろのデザートを運ばせた。\VUgras{チーズ}、\VUgras{果物}、
\VUgras{ビスケット}である。さらにサン・テミリオンのよい\VUgras{葡萄酒}を一本添えた。
食事にたいへん満足して、\VUgras{給仕} \textsuperscript{(57)} によい\VUgras{心づけ}
\textsuperscript{(59)} をやってから卓を離れた。

%%K t13-10-1 p
10. --- 私が帰りかけているのを見て、\VUgras{支配人} \textsuperscript{(60)} は
バルコニーへ上がって\VUgras{大洋} \textsuperscript{(62)} のみごとな眺めを見るように
勧めてくれた。遠くには漁の\VUgras{小舟} \textsuperscript{(63)}、\VUgras{帆船}
\textsuperscript{(64)}、それに巨大な\VUgras{郵船} \textsuperscript{(65)} が見え、その
黒い影が白い\VUgras{雲} \textsuperscript{(67)} ——\VUgras{水平線}
\textsuperscript{(66)} の雲である——の上に浮き出ていた。軽い風が\VUgras{旗}
\textsuperscript{(70)} をひるがえしていた。旗は\VUgras{看板} \textsuperscript{(69)} の
上に立ててあり、看板は\VUgras{玄関広間} \textsuperscript{(68)} のものだった。

%%K t13-tit-3 sub
\VUsaut{3.0mm}
\VUpk{10.2pt}{40}{楽しみ。}
\VUsaut{3.0mm}

%%K t13-11-1 p
11. --- 眺めがあまりに面白かったので、私は明日カフェのテラスに
\VUgras{オペラグラス} \textsuperscript{(71)} か\VUgras{望遠鏡} \textsuperscript{(72)} を
持って上がろうと決めた。

%%K t13-12-1 p
12. --- バルコニーを離れて、私はホテルのカフェへ\VUgras{飲み物}
\textsuperscript{(73)} を一杯飲みに、\VUgras{玉突きの一勝負} \textsuperscript{(75)} を
打ちに行った。立ちどまらずにカフェの広間を通り抜けた。そこでは大勢の客が
\VUgras{チェス} \textsuperscript{(78)}、\VUgras{チェッカー} \textsuperscript{(79)}、
\VUgras{ドミノ} \textsuperscript{(80)}、\VUgras{バックギャモン} \textsuperscript{(81)}、
\VUgras{トランプ} \textsuperscript{(82)} をしていた。私はまっすぐ\VUgras{玉突き場}
\textsuperscript{(74)} へ上がった。

%%K t13-13-1 p
13. --- 勝負が終わると、私は一階へおりて、給仕に\VUgras{封筒}
\textsuperscript{(84)}、\VUgras{便箋} \textsuperscript{(83)}、\VUgras{切手}
\textsuperscript{(85)}、\VUgras{ペン} \textsuperscript{(87)}、\VUgras{インク}
\textsuperscript{(86)} を頼み、手紙を書いた。

%%K t13-13-2 p
それから私は\VUgras{馭者} \textsuperscript{(92)} を呼んだ。その\VUgras{馬車}
\textsuperscript{(88)} はカフェの前に停まっていた。時間ぎめか一走りかで値段を尋ね、話が
まとまると、彼は\VUgras{車の扉} \textsuperscript{(89)} を開けて私を乗せてくれた。彼は
また\VUgras{御者台} \textsuperscript{(91)} に上がり、\VUgras{鞭} \textsuperscript{(93)}
を手早く一振りして馬を走らせ、私たちはたいそう気持ちのよい遠乗りに出かけた。
\VUsaut{6.0mm}
\VUfilet{22mm}
